\documentclass[12pt,table,xcolor={dvipsnames}]{beamer}
\usetheme{Pittsburgh}
\usecolortheme{seagull}
%\usepackage[utf8]{inputenc}
\usepackage{fontspec}
\usepackage[brazilian]{babel}
\usepackage{amsmath}
\usepackage{listings}
\usepackage{tikz}
\usetikzlibrary{calc,shapes.multipart,chains,arrows, positioning, automata}
\usepackage{multirow}
\usepackage{amsfonts}
\usepackage{amssymb}
\usepackage{lstlinebgrd}
\usepackage{graphicx}
\author{Trabalho de Implementação I \\ Simulação de um Supermercado}
\title{Estruturas de Dados}
%\setbeamercovered{transparent}
\setbeamertemplate{navigation symbols}{}
%\logo{\includegraphics[scale=0.015]{Brasao_UFSC.png}\includegraphics[scale=0.2]{brasao_PPGCC.jpg}}
\institute{Departamento de Informática e de Estatística \\ Prof. Jean Everson Martina \\ Prof. Aldo von Wangenheim}
\date{2016.2}
\subject{}
\usebackgroundtemplate{\includegraphics[width=\paperwidth,
height=\paperheight]{../reusable_images/fundo_UFSC.png}}
\begin{document}

{
\usebackgroundtemplate{\includegraphics[width=\paperwidth,
height=\paperheight]{../reusable_images/fundo_capa.png}}
\begin{frame}
\titlepage
\includegraphics[scale=0.3]{../reusable_images/brasao_INE.png}
\end{frame}
}

\begin{frame}[fragile]{Modelagem e Simulação de um Supermercado}
\lstset{language=C++,
          keywordstyle=\color{blue}\ttfamily,
          stringstyle=\color{red}\ttfamily,
          commentstyle=\color{OliveGreen}\ttfamily,
          breaklines=true,
          basicstyle=\ttfamily\footnotesize
          }
\begin{itemize}
\item Esta simulação vai permitir:
\begin{itemize}
\item Ao dono de um supermercado gerar estatísticas que permitam ver qual o melhor numero de caixas de supermercado para a freqüência de chegada de clientes do seu ponto;
\item Ao dono do supermercado realizar cálculos dos custos de operação do supermercado em função do faturamento gerado por cada caixa.
\end{itemize}
\end{itemize}
\end{frame}

\begin{frame}[fragile]{Modelo básico:}
\lstset{language=C++,
          keywordstyle=\color{blue}\ttfamily,
          stringstyle=\color{red}\ttfamily,
          commentstyle=\color{OliveGreen}\ttfamily,
          breaklines=true,
          basicstyle=\ttfamily\footnotesize
          }
\begin{itemize}
\item O modelo do supermercado terá 3 tipos de dados básicos (objetos):
\begin{itemize}
\item O Cliente - Um cliente chega no supermercado num determinado momento, compra uma quantidade de coisas, podendo pagar com cheque ou dinheiro;
\item O Caixa de Supermercado - O caixa de supermercado possui uma Fila de Clientes e características como eficiência, etc;
\item O Supermercado - Um Supermercado possui um determinado número de caixas, que são diferentes quanto à sua eficiência e um Relógio interno, além da capacidade de calcular as suas estatísticas.
\end{itemize}
\end{itemize}
\end{frame}

 \begin{frame}[fragile]{Forma básica de funcionamento do sistema}
\lstset{language=C++,
          keywordstyle=\color{blue}\ttfamily,
          stringstyle=\color{red}\ttfamily,
          commentstyle=\color{OliveGreen}\ttfamily,
          breaklines=true,
          basicstyle=\ttfamily\footnotesize
          }
\begin{itemize}
\item No sistema o usuário entra com os seguintes dados:
\begin{itemize}
\item O número de caixas do supermercado, indicando para cada caixa se ele é eficiente, médio ou ruim;
\item O tempo médio de chegada, em segundos, de novos clientes no supermercado;
\item Quanto tempo de operação, em horas, do supermercado o usuário deseja simular (tempo total de simulação);
\item Estes dados podem ser digitados em uma tela de entrada de dados ou então lidos de um arquivo de configuração. Implemente ambas as opções. O formato do arquivo de configuração está descrito abaixo.
\end{itemize}
\end{itemize}
\end{frame}


 \begin{frame}[fragile]{Forma básica de funcionamento do sistema}
\lstset{language=C++,
          keywordstyle=\color{blue}\ttfamily,
          stringstyle=\color{red}\ttfamily,
          commentstyle=\color{OliveGreen}\ttfamily,
          breaklines=true,
          basicstyle=\ttfamily\footnotesize
          }
\begin{itemize}
\item O Sistema proverá ao usuário os seguintes dados:
\begin{itemize}
\item O faturamento total do supermercado no período.
\item O faturamento médio por caixa.
\item O faturamento real de cada caixa.
\item O lucro por caixa, descontando-se o salario deste caixa no período compreendido pela simulação. Ignore outros custos, como luz elétrica, manutenção do terminal de ponto de venda, etc.
\end{itemize}
\end{itemize}
\end{frame}


 \begin{frame}[fragile]{Forma básica de funcionamento do sistema}
\lstset{language=C++,
          keywordstyle=\color{blue}\ttfamily,
          stringstyle=\color{red}\ttfamily,
          commentstyle=\color{OliveGreen}\ttfamily,
          breaklines=true,
          basicstyle=\ttfamily\footnotesize
          }
\begin{itemize}
\item O tempo médio de permanência de um cliente na fila.
\item O número de clientes que desistiram porque a fila estava muito longa.
\item O faturamento que deixou de ser realizado em função dos clientes que desistiram.
\item O Nome do Supermercado lido do arquivo de configuração deve aparecer na tela quando o sistema terminar a simulação.
\end{itemize}
\end{frame}

\begin{frame}[fragile]{Forma básica de funcionamento do Modelo}
\lstset{language=C++,
          keywordstyle=\color{blue}\ttfamily,
          stringstyle=\color{red}\ttfamily,
          commentstyle=\color{OliveGreen}\ttfamily,
          breaklines=true,
          basicstyle=\ttfamily\footnotesize
          }
\begin{itemize}
\item O supermercado funcionará da seguinte forma:
\begin{itemize}
\item A intervalos randômicos com média dada pelo usuário chegam novos clientes no supermercado.
\item Cada cliente compra um número variável de produtos, que possuem preço variável, ambos gerados aleatoriamente em um intervalo.
\item Há dois tipos de cliente: Um tipo de cliente escolhe a menor fila, o outro olha para a quantidade de compras das pessoas na fila e escolhe a com a menor quantidade para ficar.
\end{itemize}
\end{itemize}
\end{frame}

\begin{frame}[fragile]{Forma básica de funcionamento do Modelo}
\lstset{language=C++,
          keywordstyle=\color{blue}\ttfamily,
          stringstyle=\color{red}\ttfamily,
          commentstyle=\color{OliveGreen}\ttfamily,
          breaklines=true,
          basicstyle=\ttfamily\footnotesize
          }
\begin{itemize}
\item Se o cliente não encontrar uma fila que o satisfaça, ele abandona o carrinho com as compras e vai embora. O valor dos produtos de um carrinho abandonado são computados como faturamento que o supermercado deixou de realizar.
\item Se o tamanho de todas as filas ultrapassar um valor T (definido pelo usuário como parâmetro de simulação), o supermercado chamará mais um caixa para trabalhar. Conta-se o custo desse caixa em dobro (sobrehora) a partir do momento em que foi chamado. Considere o tempo que o caixa leva para ir ao supermercado como negligenciável. 
\end{itemize}
\end{frame}


\begin{frame}[fragile]{Forma básica de funcionamento do Modelo}
\lstset{language=C++,
          keywordstyle=\color{blue}\ttfamily,
          stringstyle=\color{red}\ttfamily,
          commentstyle=\color{OliveGreen}\ttfamily,
          breaklines=true,
          basicstyle=\ttfamily\footnotesize
          }
\begin{itemize}
\item O supermercado possuirá:
\begin{itemize}
\item Uma lista circular que conterá todos os caixas e que será percorrida constantemente até o fim da simulação.
\item A cada passada pela lista, o supermercado incrementa o seu relógio interno de um segundo.
\item Ao passar pela lista, toda vez que encontrar um cliente numa fila com tempo de saída igual ao tempo atual, atenderá o cliente e o tirará da fila.
\item Se o tempo atual for o tempo da chegada do próximo cliente, este será gerado, suas compras geradas e um lugar para ele numa fila escolhido. Logo após, o tempo de chegada do próximo cliente será gerado. O programa inicia com a geração da chegada do primeiro cliente.
\end{itemize}
\end{itemize}
\end{frame}

\begin{frame}[fragile]{Geração de valores aleatórios em um intervalo}
\lstset{language=C++,
          keywordstyle=\color{blue}\ttfamily,
          stringstyle=\color{red}\ttfamily,
          commentstyle=\color{OliveGreen}\ttfamily,
          breaklines=true,
          basicstyle=\ttfamily\footnotesize
          }
\begin{itemize}
\item Gerar valores aleatorios com distribuição uniforme no intervalo 0 a 1, utilize as funções rand e srand.
\item Lembre-se de inicializar o sempre gerador de numeros aleatorios, antes de usar, para garantir de que sejam usados valores diferentes em cada simulação.
\item Para gerar um numero aleatorio com distribuição uniforme em um intervalo, pegue este resultado, multiplique pelo tamanho do intervalo e adicione a valor do limite inferior do intervalo. 
\end{itemize}
\end{frame}

\begin{frame}[fragile]{Modelagem de um Cliente}
\lstset{language=C++,
          keywordstyle=\color{blue}\ttfamily,
          stringstyle=\color{red}\ttfamily,
          commentstyle=\color{OliveGreen}\ttfamily,
          breaklines=true,
          basicstyle=\ttfamily\footnotesize
          }
\begin{itemize}
\item Um cliente é uma classe que possui os seguintes atributos:
\begin{itemize}
\item Uma variavel que diz se ele busca a menor fila ou a com menos produtos;
\item A informação se ele vai pagar com cartão ou dinheiro. A probabilidade de ele pagar com cartão é de 80\%;
\item Um total de compras. Gerado entre 2 e 100;
\item Um valor total das compras. O valor de cada ítem varia entre R\$ 1 e R\$ 90. Para cada ítem o valor é gerado aleatoriamente. Depois disso o valor total é computado;
\item O tempo de chegada no supermercado;
\item O tempo em que vai sair da fila do caixa.
\end{itemize}
\end{itemize}
\end{frame}


\begin{frame}[fragile]{O tempo de saída da fila}
\lstset{language=C++,
          keywordstyle=\color{blue}\ttfamily,
          stringstyle=\color{red}\ttfamily,
          commentstyle=\color{OliveGreen}\ttfamily,
          breaklines=true,
          basicstyle=\ttfamily\footnotesize
          }
\begin{itemize}
\item Computado para o cliente depois que é escolhido um caixa para o cliente. Calculado com base nas compras do cliente, desempenho do caixa, se vai pagar com cartão e na soma dos tempos das pessoas que estão na fila
\item O sistema sempre olha o tempo de saída do primeiro da fila de cada caixa, para ver se ele vai sair agora;
\item  Se não houver fila com menos que 10 pessoas, o cliente abandona o carrinho. Assuma que o cliente que abandona o carrinho está muito irritado vai escolher outro supermercado para suas próximas duas compras, isso significa que o faturamento deixado de realizar computado para cada cliente que abandona o carrinho é multiplicado por três.
\end{itemize}
\end{frame}

\begin{frame}[fragile]{Modelagem de um Caixa}
\lstset{language=C++,
          keywordstyle=\color{blue}\ttfamily,
          stringstyle=\color{red}\ttfamily,
          commentstyle=\color{OliveGreen}\ttfamily,
          breaklines=true,
          basicstyle=\ttfamily\footnotesize
          }
\begin{itemize}
\item Um Caixa possui uma Fila de Clientes. O Caixa:
\begin{itemize}
\item Tem um identificador (Rosilda, Dona\_Beia, etc); 
\item Tem um desempenho: eficiente, médio, ruim. Um caixa eficiente leva 1 segundo para um ítem, o médio leva 2 segundos e o ruim 4 segundos. Codifique eficiente, médio, ruim como numeros inteiros 1, 2, 3.
\item Tem um salario.
\item Se o cliente paga com cartão, o caixa eficiente leva mais 10 segundos para processar, o médio leva mais 25 segundos e o ruim mais 60 segundos. Se o cliente paga em dinheiro, o tempo de pagamento é Zero.
\item O tempo de saída de um cliente é calculado quando ele é incluído na fila.
\end{itemize}
\end{itemize}
\end{frame}

\begin{frame}[fragile]{Atributos do Caixa}
\lstset{language=C++,
          keywordstyle=\color{blue}\ttfamily,
          stringstyle=\color{red}\ttfamily,
          commentstyle=\color{OliveGreen}\ttfamily,
          breaklines=true,
          basicstyle=\ttfamily\footnotesize
          }
\begin{itemize}
\item Ponteiro para a fila de clientes.
\item Número de clientes na fila.
\item Número de clientes atendidos.
\item Tempo médio de espera de um cliente na fila.
\item Faturamento total.
\item Faturamento médio.
\item Além desses, os vistos anteriormente:
\begin{itemize}
\item Identificador
\item Eficiência.
\end{itemize}
\end{itemize}
\end{frame}

\begin{frame}[fragile]{Modelagem do Supermercado}
\lstset{language=C++,
          keywordstyle=\color{blue}\ttfamily,
          stringstyle=\color{red}\ttfamily,
          commentstyle=\color{OliveGreen}\ttfamily,
          breaklines=true,
          basicstyle=\ttfamily\footnotesize
          }
\begin{itemize}
\item A classe Supermercado possui os seguintes atributos:
\begin{itemize}
\item Um ponteiro para a lista circular de Caixas.
\item Um relógio interno, inicialmente Zero.
\item O número de clientes que desistiram.
\item O valor das compras dos clientes que desistiram.
\item Todos os outros dados estatísticos e operacionais podem ser calculados a partir dos dados armazenados nos caixas.
\end{itemize}
\end{itemize}
\end{frame}

\begin{frame}[fragile]{Arquivo de Configuração}
\lstset{language=Bash,
          keywordstyle=\color{blue}\ttfamily,
          stringstyle=\color{red}\ttfamily,
          commentstyle=\color{OliveGreen}\ttfamily,
          breaklines=true,
          basicstyle=\ttfamily\footnotesize
          }

\begin{itemize}
\item Ao ser iniciado o programa, este lerá um arquivo de configuração contendo as informações necessárias para a execução da simulação. Este arquivo conterá as seguintes informações, utilizando-se da seguinte sintaxe:
\end{itemize}
\begin{lstlisting}
# comentarios quaisquer ate o fim da linha 
# mais um comentario. 
# ======================================= 
Nome do Supermercado Inclusive Brancos 
tempo_de_simulacao_em_horas 
tempo_medio_de_chegada_de_cliente_em_segundos 
numero de caixas 
id_caixa1_sem_brancos eficiencia1 salario1 
id_caixa2_sem_brancos eficiencia2 salario2 
- - - 
id_caixaN_sem_brancos eficienciaN salarioN
\end{lstlisting}
\end{frame}

\begin{frame}[fragile]{Exemplo de Arquivo de Configuração}
\lstset{language=Bash,
          keywordstyle=\color{blue}\ttfamily,
          stringstyle=\color{red}\ttfamily,
          commentstyle=\color{OliveGreen}\ttfamily,
          breaklines=true,
          basicstyle=\ttfamily\footnotesize
          }


\begin{lstlisting}
#============================= 
# Dados do Super NaoCompre 
# Programado por Hack++ 
# Arquivo: naocompre.dat 
#============================= 
NaoCompre Costeira 
5 
7 
4 
Maria_Benta 1 800 
Juliana_Digito 1 800 
Zeca_Mole 3 180 
Joao_DeMora 3 180 
#================FIM 
\end{lstlisting}
\end{frame}










{
\usebackgroundtemplate{\includegraphics[width=\paperwidth,
height=\paperheight]{../reusable_images/fundo_capa.png}}
\begin{frame}

{\LARGE Perguntas????}

\end{frame}
}


{
\usebackgroundtemplate{\includegraphics[width=\paperwidth,
height=\paperheight]{../reusable_images/fundo_capa.png}}
\begin{frame}
\includegraphics[scale=0.8]{../reusable_images/cc_logo_arge.png}\hspace{0.5cm}
\includegraphics[scale=0.95]{../reusable_images/by.png}

\vspace{1cm}
Este trabalho está licenciado sob uma Licença Creative Commons Atribuição 4.0 Internacional. Para ver uma cópia desta licença, visite http://creativecommons.org/licenses/by/4.0/.

\end{frame}
}
\end{document}
